\documentclass[11pt]{article}
\usepackage[margin=1in]{geometry}
\usepackage{amsmath,amssymb,amsthm}
\usepackage[T1]{fontenc}
\title{An infinite common periodic-point set\\
  A disproof of the finiteness assertion of conjecture 00000007757}
\author{earthking11}
\date{}

\begin{document}
\maketitle

Conjecture 00000007757 asserts that two distinct rational maps of degree at
least two in a commuting semigroup have only finitely many common periodic
points in $\mathbb P^1(\overline{\mathbb Q})$, with a cardinality bound
depending uniformly on the product of their degrees. We give a counterexample
over $\mathbb Q$.

Put
\[
  f(z)=z^2, \qquad g(z)=z^4=f^{\circ 2}(z).
\]
These are distinct rational maps of degrees $2$ and $4$. They commute because
$f(g(z))=z^8=g(f(z))$. Their polynomial extensions both fix infinity, so the
commutation holds on the entire projective line.

For each $n\geq 0$, choose a primitive $(2^{n+1}-1)$st root of unity
$\zeta_n\in\overline{\mathbb Q}$. Since
\[
  f^{\circ(n+1)}(\zeta_n)
  =\zeta_n^{2^{n+1}}
  =\zeta_n^{(2^{n+1}-1)+1}
  =\zeta_n,
\]
the point $[\zeta_n:1]$ is periodic for $f$. Also
$g^{\circ(n+1)}=f^{\circ 2(n+1)}$, so it is periodic for $g$.
The orders $2^{n+1}-1$ strictly increase with $n$. Primitive roots of
different orders are distinct, so these give infinitely many different
points of $\mathbb P^1(\overline{\mathbb Q})$ in $P_f\cap P_g$. Thus the
intersection is infinite even though $\deg f\deg g=8$ is fixed.

\paragraph{Formalization.}
The Lean file works over $K=\operatorname{AlgebraicClosure}(\mathbb Q)$.
It represents $\mathbb P^1(K)$ by its affine chart together with infinity,
using \texttt{Option K}: \texttt{some z} is $[z:1]$ and \texttt{none} is
$[1:0]$. The functions \texttt{F} and \texttt{G} extend the two polynomial
maps to this model. The file proves their commutation on all points,
distinctness, the two polynomial degrees, and the periodicity and injectivity
of the chosen roots. The final theorem
\texttt{TLMC7757.infinite\_common\_periodic} states that the actual common
periodic-point set is infinite. Its axiom audit contains only
\texttt{propext}, \texttt{Classical.choice}, and \texttt{Quot.sound}.

\paragraph{Scope.}
The example deliberately has a common iterate. The conjecture's first
finiteness assertion does not exclude this case, and its separate clause for
maps with no common iterate is not needed for the contradiction.

\end{document}
